\chapter{Procedimientos de emergencia}

En las siguientes páginas se describen instrucciones abreviadas para atender una emergencia. Procure seguir fielmente cada paso, pero tenga presente que la situación real determinará las acciones más adecuadas.

En cada procedimiento se presupone que el personal involucrado cuenta con el siguiente material:
\begin{itemize}
  \item pértiga para recuperación de fuentes;
  \item detector de radiación;
  \item detector con alarma;
  \item dosímetro de lectura directa;
  \item dosímetro termoluminiscente;
  \item blindaje de media caña;
  \item blindaje de emergencia.
\end{itemize}

\section{Falla en el mecanismo impulsor con la fuente en la manguera}

\textbf{Clasificación:} sin personas sobreexpuestas, I-A; con personas sobreexpuestas, II-A.

\begin{enumerate}
  \item[7.1.1] Acordone un área de 25 metros a la redonda de donde se encuentre la fuente.
  \item[7.1.2] Evalúe la seguridad de las personas. Observe si no hay riesgos y, si los hubiera, procure eliminarlos antes de pasar al siguiente punto.
  \item[7.1.3] Coloque el blindaje de media caña sobre la manguera, adelante del contenedor.
  \item[7.1.4] Jale el contenedor con el maneral, observando el nivel de radiación, hasta que la fuente quede cubierta con el blindaje. Estas dos maniobras deben efectuarse lo más rápidamente posible.
\end{enumerate}

\begin{figure}[htbp]
  \centering
  \includegraphics[width=0.86\textwidth]{jaleContenedor.png}
  \caption{Jalado del contenedor con el maneral mientras se observa el detector.}
  \label{fig:emergencia-jale-contenedor}
\end{figure}

\begin{figure}[htbp]
  \centering
  \includegraphics[width=0.65\textwidth]{mediaCana.png}
  \caption{Colocación del blindaje de media caña sobre la manguera.}
\end{figure}

\begin{enumerate}
  \item[7.1.5] Proceda a inspeccionar el tipo de falla.
  \item[7.1.6] Si las mangueras están estranguladas, proceda como sigue:
  \begin{enumerate}[label=\alph*)]
    \item en caso de estrangulamiento en el maneral, mida previamente el nivel de radiación en el lugar donde se efectuará la operación; con una segueta practique un corte delante del estrangulamiento y, si es necesario, incremente el blindaje de la fuente;

    \begin{figure}[htbp]
      \centering
      \includegraphics[width=0.68\textwidth]{estrangulamientoManeral.png}
      \caption{Estrangulamiento en el maneral.}
      \label{fig:emergencia-estrangulamiento-maneral}
    \end{figure}

    \item jale rápidamente el cable de la fuente para que esta se introduzca dentro de su blindaje;

    \begin{figure}[htbp]
      \centering
      \includegraphics[width=0.86\textwidth]{jaleCable.png}
      \caption{Jalado rápido del cable de la fuente para introducirla en el blindaje.}
      \label{fig:emergencia-jale-cable}
    \end{figure}

    \item en caso de estrangulamiento en la manguera posicionadora, consulte el punto~7.4.

    \begin{figure}[htbp]
      \centering
      \includegraphics[width=0.62\textwidth]{estrangulamiento_posicionadora.png}
      \caption{Estrangulamiento en la manguera posicionadora.}
      \label{fig:emergencia-estrangulamiento-posicionadora}
    \end{figure}
  \end{enumerate}
  \item[7.1.7] Si el problema se encuentra en la pistola, desármela y proceda igualmente a jalar el cable hasta que la fuente entre en su blindaje.
  \item[7.1.8] Verifique su dosímetro para anotar la dosis que recibió durante la maniobra.
  \item[7.1.9] Reporte el incidente al Departamento de Seguridad Radiológica.
\end{enumerate}

\section{Daño en el candado del contenedor con la fuente en la manguera}

\textbf{Clasificación:} sin personas sobreexpuestas, I-A; con personas sobreexpuestas, II-A.

\begin{enumerate}
  \item[7.2.1] Acordone un área de 25 metros a la redonda de donde se encuentre la fuente.
  \item[7.2.2] Evalúe la seguridad de las personas. Observe si no hay riesgos y, si los hubiera, procure eliminarlos antes de pasar al siguiente punto.
  \item[7.2.3] Introduzca la fuente dentro del contenedor.
  \item[7.2.4] Haga que su compañero desensamble la manguera posicionadora y coloque el tapón de salida mientras usted detiene la pistola y evita que la fuente pueda moverse de su posición segura.
  \item[7.2.5] Desensamble el maneral y coloque el tapón de entrada.
  \item[7.2.6] Consulte su dosímetro y anote su dosis.
  \item[7.2.7] Notifique al Departamento de Seguridad Radiológica el incidente.
\end{enumerate}

\section{Caída del contenedor}

\textbf{Clasificación:} I-A.

\begin{enumerate}
  \item[7.3.1] Defina un área restringida de 25 metros alrededor de la fuente y acordónela.
  \item[7.3.2] Introduzca la fuente a su contenedor.
  \item[7.3.3] Si no es posible, observe los daños en las mangueras, específicamente en las conexiones con el contenedor.
  \item[7.3.4] Coloque el blindaje de media caña delante del contenedor.
  \item[7.3.5] Jale el contenedor con el maneral, observando el nivel de radiación, hasta que la fuente quede cubierta con el blindaje. Estas dos maniobras deben efectuarse lo más rápidamente posible.

  \begin{figure}[htbp]
    \centering
    \includegraphics[width=0.86\textwidth]{jaleContenedor.png}
    \caption{Jalado del contenedor durante la atención de una caída.}
    \label{fig:emergencia-caida-jale}
  \end{figure}

  \item[7.3.6] Inspeccione el daño y vea si puede arreglarlo. El objetivo es llevar la fuente hacia dentro de su contenedor.
  \item[7.3.7] En caso de no ser posible, habrá que romper la manguera y soltar la fuente al piso.
  \item[7.3.8] Para este caso, consulte el procedimiento «Fuente atorada en la manguera».
  \item[7.3.9] Si la fuente se encuentra en su contenedor, inspeccione cuidadosamente los niveles de radiación mínimos en todas las direcciones. Verifique las conexiones y el funcionamiento del candado.
  \item[7.3.10] En caso de altos niveles de radiación (200\,mR/h), interrumpa su uso y reporte el incidente al Departamento de Seguridad Radiológica.
  \item[7.3.11] Hasta en tanto no reciba instrucciones, transporte el contenedor a una fosa de seguridad o al almacén de material radiactivo, empleando una plataforma rodante o cualquier medio que permita la manipulación remota del contenedor.
  \item[7.3.12] El Departamento de Seguridad Radiológica practicará una prueba de fuga a la brevedad posible e inspeccionará los daños causados al equipo.
\end{enumerate}

\section{Fuente atorada en la manguera}

\textbf{Clasificación:} I-A.
\begin{enumerate}
  \item[7.4.1] Acordone un área de 25 metros a la redonda de donde se encuentra la fuente.
  \item[7.4.2] Consulte su dosímetro y póngalo en cero antes de iniciar las maniobras. Anote la dosis acumulada.
  \item[7.4.3] Coloque el blindaje de media caña adelante del contenedor, sobre la manguera. Efectúe esta maniobra rápidamente.
  \item[7.4.4] Tomando la mayor distancia que le sea posible, jale con el maneral al contenedor hasta que se observe en el detector que la fuente ha quedado blindada.
\end{enumerate}

\begin{figure}[htbp]
  \centering
  \includegraphics[width=0.88\textwidth]{jale_maneral_fuente.png}
  \caption{Jalado del contenedor desde la mayor distancia posible.}
  \label{fig:emergencia-fuente-atorada}
\end{figure}

Planee su recorrido de tal manera que, de caerse durante esa maniobra, la fuente no quede en una posición difícil de volver a sujetar.

\textit{Nota de revisión:} en las páginas fuente disponibles, la numeración salta de 7.4.4 a 7.4.12. Los pasos 7.4.5--7.4.11 no se reconstruyen aquí porque no aparecen en las imágenes revisadas; no se ha inventado texto para completar el salto.

\begin{enumerate}
  \item[7.4.12] Asegure la fuente y llévela al almacén de material radiactivo.
  \item[7.4.13] Con otro maneral, conecte la fuente e introdúzcala al contenedor.
\end{enumerate}

\section{Inundación del almacén de material radiactivo}

\textbf{Clasificación:} I-A.
\begin{enumerate}
  \item[7.5.1] Efectúe una inspección visual verificando la integridad de los accesos al almacén.
  \item[7.5.2] Si hay indicio de que algunos objetos hayan podido ser arrastrados por las aguas, verifique a la brevedad posible el inventario de fuentes y efectúe las maniobras necesarias para este efecto.
  \item[7.5.3] Si todo el material se encuentra presente, haga un levantamiento de niveles de radiación e inspeccione una a una todas las fuentes.
  \item[7.5.4] Proceda a una limpieza detallada y lubricación del equipo.
  \item[7.5.5] En caso de que falte alguna fuente, proceda a su búsqueda inmediatamente, verificando en primer término los lugares cercanos. Observe las huellas que el agua deja a su paso y, donde se encuentren obstáculos naturales, inspecciónelos detenidamente.
  \item[7.5.6] Notifique inmediatamente al Departamento de Seguridad Radiológica y a la C.N.S.N.S.
  \item[7.5.7] Declare una emergencia tipo II-B y proceda a la formación de brigadas de búsqueda, así como a asignar las responsabilidades de comunicación a alguna persona.
  \item[7.5.8] Proceda como en el punto 7.12.
\end{enumerate}

\section{Derrumbe de construcción o almacén de material radiactivo}

\textbf{Clasificación:} I-A.
\begin{enumerate}
  \item[7.6.1] Efectúe una inspección con el detector y observe si el nivel de radiación es seguro para efectuar las maniobras de remoción de escombros.
  \item[7.6.2] Durante las maniobras vigile constantemente el nivel de radiación y asegúrese de portar sus dosímetros de lectura directa y termoluminiscente.
  \item[7.6.3] Si hay sospechas de que alguna fuente se encuentre descubierta, haga que las maniobras se efectúen con maquinaria para que haya un factor de protección suficiente para el personal.
  \item[7.6.4] Dé instrucciones para que la maquinaria separe porciones de escombro hasta localizar la o las fuentes causantes de la alta radiación.
  \item[7.6.5] Proceda a su recuperación utilizando herramientas remotas o pértiga.
  \item[7.6.6] Coloque las fuentes en blindajes de emergencia.
  \item[7.6.7] Proceda a efectuar pruebas de fuga a todas las fuentes que se encuentren en el lugar del incidente.
\end{enumerate}

\section{Incendio o explosión en almacén de material radiactivo}

\textbf{Clasificación:} I-A.
\begin{enumerate}
  \item[7.7.1] Efectúe una verificación con su detector. Asegúrese de portar sus dosímetros de lectura directa y termoluminiscentes.
  \item[7.7.2] Si los niveles de radiación son seguros, proceda a inspeccionar en detalle los contenedores de fuentes. Reporte sus resultados al Departamento de Seguridad Radiológica.
  \item[7.7.3] Practique pruebas de fuga a todas las fuentes implicadas.
  \item[7.7.4] Si llegara a encontrar contaminación removible con las pruebas de frotis, retírese del lugar inspeccionando sus ropas y calzado, acordone el área y evite la remoción de objetos del área restringida.
  \item[7.7.5] Anote nombres y direcciones de personas que pudieran haber resultado afectadas por la explosión o contaminación e informe a ambulancias y hospitales la presencia de radiación y posible contaminación de los pacientes.
  \item[7.7.6] Evite fumar, ingerir alimentos o bebidas cerca del área del incidente. Notifique a la C.N.S.N.S.
  \item[7.7.7] Si no encontrara contaminación pero los niveles de radiación fueran altos, observe si puede descubrir su procedencia.
  \item[7.7.8] Con cuidado, vaya aislando las diferentes fuentes, alejándolas del lugar hasta encontrar la o las causantes de los niveles de radiación.
  \item[7.7.9] Si la procedencia es de algún contenedor dañado, prepare un baño grande con agua de aproximadamente 35--45\,cm de profundidad.
  \item[7.7.10] Vea la manera de introducir el contenedor en el baño. Esto disminuirá el nivel de radiación.
  \item[7.7.11] Vea si puede retirar los tapones del equipo. Utilice las herramientas que sean necesarias.
  \item[7.7.12] Una vez hecho esto, conecte la fuente a un maneral y vea si puede retirarla del contenedor dañado.
  \item[7.7.13] Si esto es posible, trasládela a un contenedor seguro.
  \item[7.7.14] Si la fuente no puede ser retirada del equipo, abandone la tarea y prepare un blindaje adecuado para todo el contenedor. Trasládelo como desecho al ININ o, si es posible, al proveedor.
\end{enumerate}

\section{Caída de la fuente en ríos, pantanos o mar}

\textbf{Clasificación:} I-B.
\begin{enumerate}
  \item[7.8.1] Defina un área restringida de acuerdo con las posibilidades del caso.
  \item[7.8.2] Si le es posible, marque con una boya el punto aproximado donde se sospecha que cayó la fuente. Puede funcionar como boya algún material flotante amarrado a una cuerda larga y luego a una pica. Puede usar también un pequeño dibujo tomando como referencia objetos situados en ambas orillas, por ejemplo, de un río.
  \item[7.8.3] Notifique al Departamento de Seguridad Radiológica y a la C.N.S.N.S.
  \item[7.8.4] Prepare un plan de rescate y consúltelo con el Grupo de Soporte; no intente maniobras arriesgadas sin antes obtener su aprobación o la de la C.N.S.N.S.
  \item[7.8.5] Analice la conveniencia de obtener el apoyo de las autoridades y de notificar al público.
  \item[7.8.6] Defina el curso de las acciones de común acuerdo con la C.N.S.N.S.
\end{enumerate}

\section{Fuente suelta en el piso}

\textbf{Clasificación:} sin personal sobreexpuesto, I-B; con personal sobreexpuesto, II-B.
\begin{enumerate}
  \item[7.9.1] Defina un área restringida de 25 metros a la redonda de donde se encuentra la fuente.
  \item[7.9.2] Recargue sus dosímetros de lectura directa, anotando su dosis.
  \item[7.9.3] Coloque el blindaje de emergencia en un área abierta y libre de obstáculos. Coloque un embudo en la boquilla superior.
  \item[7.9.4] Con la pértiga sujete la fuente y cuidadosamente llévela hasta posicionarla en el embudo del contenedor.

  \begin{figure}[htbp]
    \centering
    \includegraphics[width=0.88\textwidth]{recuperacion_pertiga.png}
    \caption{Introducción de la fuente en el contenedor mediante una pértiga.}
    \label{fig:emergencia-recuperacion-pertiga}
  \end{figure}

  \item[7.9.5] Planee su recorrido de tal manera que, de caerse durante esa maniobra, la fuente no quede en posición difícil de volver a sujetar. Efectúe esta maniobra rápidamente.
  \item[7.9.6] Una vez en el blindaje, la fuente se encuentra segura.
  \item[7.9.7] Desconecte la manguera posicionadora del contenedor y haga que el cable impulsor de la fuente pase por el contenedor hasta conectarlo a la fuente.
  \item[7.9.8] Retírese hasta donde el maneral le permita y, accionando la pistola, retraiga la fuente hasta su contenedor original. Verifique la dosis que recibió y anótela.

  \begin{figure}[htbp]
    \centering
    \includegraphics[width=0.88\textwidth]{retraccion_fuente.png}
    \caption{Retracción de la fuente hacia el contenedor original.}
    \label{fig:emergencia-retraccion-fuente}
  \end{figure}

  \item[7.9.9] En caso de que su dosímetro haya resultado sobreexpuesto, haga una relación cuidadosa y detallada de todas las actividades realizadas e inclúyala en el reporte de emergencias radiológicas. Mencione los tiempos empleados en cada maniobra y las distancias a las que se mantuvo de la fuente. El Departamento de Seguridad Radiológica calculará su dosis con esta información y enviará a leer los dosímetros TLD inmediatamente.
  \item[7.9.10] Practíquese análisis bioquímicos a la brevedad posible (antes de 24 horas) que incluyan biometría hemática, recuento de plaquetas y espermatobioscopia. Envíelos al Departamento de Seguridad Radiológica y espere instrucciones.
\end{enumerate}

\section{Sobreexposición de personal o público}

\textbf{Clasificación:} II-A.
\begin{enumerate}
  \item[7.10.1] Si la sobreexposición se sospecha porque uno de los dosímetros de lectura directa se encuentra saturado y no existe otro indicio de anormalidad, notifique al Departamento de Seguridad Radiológica y proceda a leer el dosímetro TLD.
  \item[7.10.2] Si los dosímetros tanto del auxiliar como del técnico se encuentran saturados, si algunas radiografías han sido veladas o si se encuentra alguna anormalidad que indique que la sobreexposición es real, proceda como sigue.
  \item[7.10.3] Vea que la fuente se encuentra segura en su contenedor.
  \item[7.10.4] Haga una revisión mental de todas las maniobras ejecutadas y anótelas, especificando: dosis que registraba el dosímetro al inicio de la jornada o la última vez que lo consultó; número y tipo de maniobras; tiempo aproximado de cada una; distancia al contenedor y las mangueras; y actividad de la fuente.
  \item[7.10.5] Notifique al Departamento de Seguridad Radiológica estos datos y cualquier otro que se juzgue importante, por ejemplo, el número de placas veladas.
  \item[7.10.6] Practíquese los siguientes análisis bioquímicos: biometría hemática, espermatobioscopia y recuento de plaquetas.
  \item[7.10.7] No se exponga a ninguna fuente de radiación ionizante hasta tener una estimación de la dosis recibida.
  \item[7.10.8] En caso de que personas del público hayan resultado sobreexpuestas, proceda de manera similar a 7.10.6 y 7.10.7.
  \item[7.10.9] Consulte con la C.N.S.N.S. en el curso de las siguientes acciones.
\end{enumerate}

\section{Accidente automovilístico}

\textbf{Clasificación:} II-B.
\begin{enumerate}
  \item[7.11.1] Evalúe la seguridad física del personal.
  \item[7.11.2] De ser necesaria alguna acción para atender a personas heridas, proceda de inmediato.
  \item[7.11.3] Evalúe la integridad de la fuente haciendo un levantamiento de niveles de radiación.
  \item[7.11.4] Si no encuentra anormalidades, lleve la fuente hasta algún lugar seguro.
  \item[7.11.5] Si encuentra alta radiación, delimite un área restringida. Notifique y consiga la cooperación de las autoridades de tránsito o de la Federal de Caminos, si se encuentran presentes. Retire a las personas lesionadas a un lugar seguro, reduciendo al mínimo el contacto personal directo.
  \item[7.11.6] Impida la entrada del público al área de emergencia y evite que saque objetos del área afectada. Advierta al público sobre el material radiactivo.
  \item[7.11.7] Anote nombres y direcciones de las personas que se encontraban cerca del área del incidente y verifique si alguna pudo haber resultado sobreexpuesta o contaminada.
  \item[7.11.8] Notifique a las oficinas generales la declaración de una emergencia.
  \item[7.11.9] Trate de determinar si existe contaminación mediante pruebas de frotis tomadas remotamente.
  \item[7.11.10] En caso de encontrar contaminación, acordone el área y determine las maniobras a seguir de acuerdo con el Grupo de Soporte o la C.N.S.N.S. Informe a ambulancias, hospitales y personal de emergencia la presencia de radiación y contaminación de los pacientes.
  \item[7.11.11] Evite la manipulación de los materiales que pudieran estar contaminados y absténgase de fumar, comer o beber cerca del área afectada.
  \item[7.11.12] Una vez que la fuente se encuentre segura, practíquele una prueba de fuga y repórtela a la C.N.S.N.S.
  \item[7.11.13] Espere instrucciones por parte del Grupo de Soporte.
\end{enumerate}

\section{Pérdida de la fuente radiactiva}

\textbf{Clasificación:} II-B.
\begin{enumerate}
  \item[7.12.1] Si la fuente ha sido robada, notifique inmediatamente a las oficinas generales y a la C.N.S.N.S.
  \item[7.12.2] Forme brigadas de búsqueda y establezca un Centro de Control de Emergencias (CCE) desde donde se puedan coordinar las acciones y centralizar la información.
  \item[7.12.3] Recabe la información necesaria para planear las labores de búsqueda (planos de la ciudad, personal disponible, vehículos, etcétera) de manera sistemática y de tal forma que se opere las 24 horas.
  \item[7.12.4] Coloque avisos para el público en los lugares aledaños donde se haya extraviado la fuente radiactiva, previniendo del peligro que representa. Incluya un dibujo o fotografía y advierta que no deben tocarla, que deben lavarse las manos y avisar de inmediato en caso de encontrarla.
  \item[7.12.5] Coordine sus acciones con las autoridades civiles y de común acuerdo con la C.N.S.N.S.
  \item[7.12.6] Si la fuente se perdió durante las maniobras de trabajo normal, forme brigadas de búsqueda y rastree sistemáticamente el área de trabajo.
  \item[7.12.7] Coloque su alarma en posición de alta sensibilidad, su detector en la escala más baja e inicie la búsqueda a la brevedad posible.
  \item[7.12.8] Evacue el área de trabajo.
  \item[7.12.9] Notifique a las oficinas generales y a la C.N.S.N.S.
  \item[7.12.10] Revise las placas tomadas para determinar si la fuente se perdió desde el inicio de la jornada de trabajo; esto le ayudará a precisar el lugar más probable de su pérdida.
  \item[7.12.11] Tan pronto la localice, proceda a su recuperación y evalúe las posibles personas que pudieron estar expuestas durante el tiempo que permaneció perdida.
  \item[7.12.12] Obtenga un informe detallado de los tiempos y tipos de maniobras que ejecutó cada una de las personas con sospecha de sobreexposición y actúe conforme al punto 7.10.
\end{enumerate}
